\section{Method}
\label{sec:overview}
\label{sec:method}

\begin{figure*}[t]
\centering
\includegraphics[width=\textwidth]{Figures_1.png}
\caption{Overview of \tool. The system extracts field-level rules from a
standard, aligns them with implementation evidence, and uses bounded reasoning
and validation tests to produce reports and reproduction artifacts.}
\label{fig:overview}
\end{figure*}

Figure~\ref{fig:overview} shows the workflow of \tool. Given a protocol
standard $S$ and an implementation $P$, \tool{} checks whether $P$ satisfies
the normative field-level requirements stated in $S$. The system first builds a
field space from the standard, extracts structured rules over those fields,
maps each rule to implementation variables and code contexts, and finally uses
an evidence-bounded reasoning-and-test agent to decide whether the observed
implementation behavior is consistent with the requirement. The output is not
only a label, but also the standard evidence, code evidence, uncertainty notes,
validation results when available, and reproduction artifacts for maintainer
triage.

\subsection{Problem Definition}

\tool{} represents each standard-derived requirement as a rule
$r=\langle f,C,A,M,E\rangle$. Here, $f$ is the constrained protocol field, $C$
is the condition under which the rule applies, $A$ is the required action or
constraint, $M$ is the normative modality such as \emph{MUST} or
\emph{SHOULD}, and $E$ records source evidence from the standard. For each
rule, the checker returns
$o=\langle y,E,U,T,R\rangle$, where
$y\in\{\textsc{Consistent},\textsc{Inconsistent},\textsc{Insufficient}\}$ is
the judgment, $E$ contains both document and code evidence, $U$ records missing
or uncertain context, $T$ stores validation-test evidence when the validation
stage runs, and $R$ is an optional inconsistency reason used only when
$y=\textsc{Inconsistent}$.

The representation separates standard-derived obligations from later
implementation judgments. A rule records what the standard requires and where
that requirement appears; a judgment records whether the implementation
evidence satisfies it. \tool{} therefore does not rely on naming similarity
alone. It extracts a structured intermediate representation from the standard
and maps that representation to variables, code contexts, and runtime behavior
using both semantic and syntactic evidence.

\subsection{Field-Space Construction}
\label{sec:fieldspace}

Directly extracting requirements from open-ended prose is unstable because
standards mix normative rules with design motivation, compatibility advice,
examples, and implementation guidance. \tool{} therefore starts by constructing
a \emph{field space}: a set of protocol objects that later rules may reference.
Field candidates are collected from message formats, extension structures,
parameter tables, and field-definition paragraphs. For each field, \tool{}
records its canonical name, aliases, parent message or structure, type,
allowed values when available, and source locations.

\textbf{Document preprocessing.}
\tool{} ingests RFC-style text, PDF, DOC, and TXT documents. To preserve local
semantic coherence, it splits the document into paragraphs and greedily merges
consecutive paragraphs into bounded text blocks without splitting a paragraph
in the middle. When adding the next paragraph would exceed the configured
threshold, the current block is committed and a new block starts from that
paragraph. The default threshold is 6000 characters and can be adjusted by
configuration. These stable, tokenizer-independent blocks are used both for
field discovery and for later rule extraction.

The field space gives the extractor a closed set of referents. This reduces the
chance that descriptive notions such as ``compatibility mode'' are mistaken
for checkable fields, while still preserving cross-document identity: the same
field may appear as a full path, short name, structure member, or contextual
pronoun. \tool{} treats two mentions as the same field only when their parent
context and semantic role are compatible. Generic names such as \emph{length},
\emph{type}, or \emph{flags} are therefore not merged across unrelated
messages. During rule extraction, every field reference must resolve to one or
more field-space entries; ambiguous references are kept as uncertainty rather
than silently collapsed. Entries with missing parent context, incompatible
aliases, or only weak prose evidence are marked ambiguous or out of scope. This
conservative policy favors precision over recall because an incorrect field
merge can mislead every downstream alignment and reasoning step.

\subsection{Field-Centric Rule Extraction}
\label{sec:rules}

After constructing the field space, \tool{} extracts field-centric rules from
normative standard sections. For each section, the model receives the section
text, the field-space entries relevant to that section, and a lexicon of RFC
2119-style normative keywords. Within this constrained context, it identifies
normative statements, resolves mentioned protocol objects to field-space
entries, and emits records under the schema in Table~\ref{tab:ruleschema}.
\emph{MAY}-level statements are filtered out because they do not impose
required behavior in our checking scope.

For example, RFC 8446 defines the TLS 1.3 \texttt{Cookie} extension body as
\texttt{opaque cookie<1..2\^{}16-1>}, so a zero-length cookie in a
\texttt{HelloRetryRequest} is invalid~\cite{rfc8446}. At this stage, \tool{}
has only analyzed the standard and extracted a rule; the reasoning-and-test
agent has not yet judged any implementation behavior:
\begin{lstlisting}[style=jsonsnippet]
{
  "f": "Cookie.cookie",
  "C": "Client receives a TLS 1.3 HRR Cookie.",
  "A": "Reject if Cookie is empty.",
  "M": "MUST",
  "E": "RFC8446 Sec.4.2.2: cookie<1..2^16-1>"
}
\end{lstlisting}

\begin{table}[t]
\centering
\caption{Rule schema used by \tool.}
\label{tab:ruleschema}
\small
\begin{tabularx}{\columnwidth}{@{}l>{\raggedright\arraybackslash}X@{}}
\toprule
\textbf{Attribute} & \textbf{Description} \\
\midrule
Field & Constrained protocol field or parameter \\
Condition & Message type, state, phase, or field predicate \\
Action & Presence, value, derivation, relation, or rejection \\
Modality & MUST, MUST NOT, SHOULD, or equivalent phrase \\
Evidence & Source section, sentence span, and referenced text \\
\bottomrule
\end{tabularx}
\end{table}

\textbf{Local scope and validation.}
This design confines rule extraction to local sections and a known field set,
instead of asking the model to infer rules from the whole standard at once.
The field space provides a traceable referential basis for protocol objects, so
each emitted rule can be linked back to concrete fields and source text. For
normative sentences that contain multiple constraints, \tool{} splits them into
multiple field-level rules, each corresponding to an independent checkable
behavior such as field presence, value constraints, range checks, state
updates, or rejection behavior.

\textbf{Normalization and validation.}
The action families are close to implementation behavior: \emph{assign},
\emph{derive}, \emph{select}, \emph{require-present},
\emph{require-absent}, \emph{compare}, \emph{range-check},
\emph{state-update}, and \emph{reject}. When one sentence contains multiple
constraints, \tool{} atomizes it into multiple rules so each rule can be
aligned and tested independently. Negated requirements are rewritten into
explicit actions, e.g., ``MUST NOT accept $x$'' becomes an expected
\emph{reject} action under condition $x$. Cross-references are expanded when
the referenced section defines a relevant field or bound, while the normalized
rule remains linked to the original text.

A deterministic checker validates the model output before alignment. It checks
that field references resolve to field-space entries, the cited source span
exists, the action type belongs to a supported action class, and the candidate
does not duplicate an existing record. Invalid outputs are retried with a
concise repair prompt; candidates that still fail are discarded rather than
turned into unauditable rules. Normative text outside the field-level scope,
such as performance guidance or cryptographic assumptions requiring proof, is
recorded as out of scope instead of silently lost.

\subsection{Syntax-Aware Rule-to-Code Alignment}
\label{sec:alignment}

The alignment stage maps each rule $r$ to implementation variables and code
contexts. It combines semantic ranking with syntax-aware evidence because pure
semantic retrieval can select comments that resemble the standard but do not
manipulate the constrained field, while pure exact matching misses renamed or
split representations.

\textbf{File ranking and variable localization.}
\tool{} first summarizes candidate C/C++ files using paths, function
signatures, structure and enum definitions, comments, and hits on field
aliases. An LLM-based ranker uses the rule, action class, aliases, and file
summaries to select files likely to contain the field representation or its
checking logic. For the constrained field, \tool{} constructs aliases by
tokenizing field names, normalizing naming styles, and expanding common
abbreviations, parent message names, and parent structure names. Within the
ranked files, \tool{} builds a Tree-sitter-based symbol index~\cite{tree_sitter}
covering variables, structure members, enum constants, macros, and function
parameters. For each indexed entity, it records declaration locations, use
locations, scope, and enclosing function. Candidate identifiers are then ranked
by alias matches, token overlap, scope, and syntactic usage patterns. A field
candidate receives higher priority when it appears in assignments,
comparisons, bounds checks, state updates, or error-return logic relevant to
the rule action, and lower priority when it is a generic temporary or lacks
action-relevant use.

\textbf{Context retrieval.}
For each candidate variable $v$, \tool{} retrieves snippets from exact textual
hits, enclosing symbol scopes, and semantic windows whose tokens overlap the
rule. Before LLM filtering, each snippet $c$ receives the score
\begin{equation}
\begin{aligned}
S_{\mathrm{ctx}}(c)
&=
 w_{\mathrm{exact}}\cdot \mathbb{I}_{\mathrm{exact}}(c)
 +w_{\mathrm{var}}\cdot |\mathcal{T}_c \cap \mathcal{T}_v|\\
&\quad
 +w_{\mathrm{rule}}\cdot |\mathcal{T}_c \cap \mathcal{T}_r|
 +b(c).
\end{aligned}
\end{equation}
Here, $\mathbb{I}_{\mathrm{exact}}$ denotes exact variable or alias occurrence,
$\mathcal{T}_c$, $\mathcal{T}_v$, and $\mathcal{T}_r$ are the code, variable,
and rule token sets, and $b(c)$ rewards action-relevant usage such as
initialization, validation, comparison, rejection, error propagation,
configuration reads, and state updates. We use
$w_{\mathrm{exact}}=5$, $w_{\mathrm{var}}=0.9$, and
$w_{\mathrm{rule}}=0.28$ as fixed defaults. These constants encode an ordinal
retrieval preference: direct field evidence should dominate variable-token
overlap, which should dominate weak rule-token overlap. Snippets without exact
variable or alias occurrence are capped to prevent broad semantic similarity
from overwhelming direct implementation evidence.

The top snippets are then sent to an LLM semantic filter. The filter sees only
the rule and ranked snippets, selects the contexts most relevant to the rule,
and explains the correspondence between each selected snippet and the expected
behavior. The result is an evidence bundle containing field representations,
checking logic, call relationships, configuration conditions, and error paths
for the next stage.

\subsection{Consistency Reasoning and Validation}
\label{sec:reasoning}
\label{sec:testagent}

After rule-to-code alignment, the reasoning-and-test agent receives an evidence
bundle containing the field rule, standard span, localized variables, selected
code snippets, call relationships, and configuration conditions. It first
performs static consistency reasoning over this bounded bundle. If the evidence
implements the required check or handling, the label is
\textsc{Consistent}; if it contradicts the rule, the label is
\textsc{Inconsistent}; if a key path is missing, such as a caller, macro
definition, configuration branch, validation helper, or error-propagation path,
the label is \textsc{Insufficient}. In the last case, $U$ records both the
uncertainty and the concrete context request for the next retrieval round.

\begin{algorithm}[!b]
\small
\DontPrintSemicolon
\KwIn{Rule $r$, implementation index $I$, maximum rounds $K$}
\KwOut{Consistency tuple $\langle y,E,U,T,R\rangle$}
$C_0 \leftarrow$ AlignRuleToCode$(r,I)$\;
\For{$i \leftarrow 0$ \KwTo $K$}{
  $o_i \leftarrow$ Reason$(r,C_i)$\;
  \If{$o_i.y \neq \mathrm{Insufficient}$}{
    \Return $o_i$\;
  }
  $q_i \leftarrow$ ContextRequest$(o_i.U)$\;
  $C_{i+1} \leftarrow C_i \cup$ RetrieveAdditionalContext$(q_i,I)$\;
  \If{$C_{i+1}=C_i$}{
    \Return $o_i$\;
  }
}
\Return $o_K$\;
\caption{Evidence-driven consistency reasoning.}
\label{alg:reasoning}
\end{algorithm}

Algorithm~\ref{alg:reasoning} summarizes this evidence-completion loop. A
determinate label stops the loop. An \textsc{Insufficient} result instead
turns the missing context into a targeted retrieval query, and the returned
context is merged into the next reasoning round. This explicit insufficient
state prevents \tool{} from converting missing evidence into a false
inconsistency while still making progress toward the code path that can settle
the rule.

Static reasoning alone can still be wrong: a check absent from a local snippet
may be enforced by an upper-level caller, while build options and conditional
compilation can change the active path. For candidate inconsistencies, the
validation-test agent therefore receives the rule, defect hypothesis, and code
evidence, then generates a rule-directed test plan. It attempts to construct
inputs that violate the standard requirement but are accepted by the
implementation, or inputs that satisfy the requirement but are rejected. The
runner can extend an existing unit test, build a small harness around a parser
or validation function, or drive a command-line implementation with generated
inputs.

The OpenSSL empty-cookie case follows this sequence. The extracted rule says
that a TLS 1.3 client must reject an empty \texttt{Cookie}. Static reasoning
then checks the aligned parser context and finds that the client reads and
stores the cookie using \texttt{PACKET\_as\_length\_prefixed\_2} and
\texttt{PACKET\_memdup}, while the reported snippet has no explicit
zero-length rejection. This is a defect hypothesis. The validation test then
injects an empty \texttt{HelloRetryRequest} cookie and observes that the client
sends a second \texttt{ClientHello} before failing later~\cite{opensslissue30868}.
The final report object combines the rule, code evidence, runtime observation,
and inconsistency reason:
\begin{lstlisting}[style=jsonsnippet]
{
  "y": "Inconsistent",
  "E": {
    "document": {
      "section": "RFC8446 Sec.4.2.2",
      "text": "opaque cookie<1..2^16-1>"
    },
    "code": {
      "file": "ssl/statem/extensions_clnt.c",
      "lines": "2372-2377",
      "text": "PACKET_as_length_prefixed_2(...) || PACKET_memdup(...); no len==0 check"
    }
  },
  "U": [],
  "T": {
    "input": "HRR with an empty Cookie extension",
    "oracle": "Client should abort immediately",
    "observed": "Client sends a second ClientHello before later failure"
  },
  "R": "When a TLS 1.3 HelloRetryRequest carries an empty Cookie, OpenSSL does not strictly enforce RFC 8446. According to the RFC,..."
}
\end{lstlisting}

\tool{} confirms an inconsistency only when the validation run reaches the
target path and the observed behavior contradicts the field rule. For confirmed
inconsistencies, the system generates a structured report and reproduction
artifacts, recording the field rule, related standard text, key code snippets,
inconsistency rationale, test input, build configuration, execution command,
observed result, and log summary. This forms a traceable evidence chain from
the standard requirement to code evidence and then to runtime behavior,
supporting later review and reproduction. For unconfirmed results, the system
does not promote the static defect hypothesis to a confirmed issue; instead, it
records the blocking condition, such as a build failure, an unreachable target
path, a missing test entry point, or an insufficient observation signal.
